\documentclass[a4paper]{article}
\usepackage[margin=1.15cm,footskip=0pt,headsep=0pt]{geometry}
\usepackage{amsmath,amssymb}
\usepackage{xcolor}
\usepackage{tikz}
\usetikzlibrary{arrows.meta,positioning,shapes.geometric}
\usepackage{hyperref}
\usepackage{fontspec}

\makeatletter
\newcommand{\MathCollectFalse}{}
\newcommand{\MathCollectTrue}{}
\let\@kernel@tabular@init\relax
\makeatother

\usepackage{xepersian}

\settextfont[
  Path = Peyda/,
  BoldFont = Peyda-Bold.ttf,
  ItalicFont = Peyda-Medium.ttf,
  BoldItalicFont = Peyda-Bold.ttf
]{Peyda-Regular.ttf}

\setlatintextfont{Times New Roman}

\IfFileExists{Peyda/Symbola.ttf}{%
  \newfontfamily\symb[Path=Peyda/]{Symbola.ttf}%
}{%
  \IfFileExists{editor/Peyda/Symbola.ttf}{%
    \newfontfamily\symb[Path=editor/Peyda/]{Symbola.ttf}%
  }{%
    \newfontfamily\symb{Symbola}%
  }%
}
\newcommand{\emoji}[1]{{\symb #1}}

\linespread{1.12}
\pagestyle{empty}
\thispagestyle{empty}

\raggedbottom
\setlength{\parindent}{0pt}
\setlength{\parskip}{0pt}

\definecolor{hlgray}{HTML}{E2E8F0}
\definecolor{hlyellow}{HTML}{FEF08A}
\definecolor{hlred}{HTML}{FEE2E2}
\definecolor{darkred}{HTML}{991B1B}
\definecolor{navy}{HTML}{0A2540}
\definecolor{teal}{HTML}{0D9488}

% =========================================================================
% PLATFORM TAXONOMY ROUTING (patho.dentnet.space)
% =========================================================================
% DOC_ID: capillary_hemangioma
% DOC_SLUG: capillary_hemangioma
% DOC_CATEGORY: cat_benign_mesenchymal
% DOC_TITLE: همانژیوم کاپیلاری (Capillary Hemangioma)
% DOC_ALIASES: Capillary Hemangioma, Infantile Hemangioma, Vascular Anomalies, همانژیوم کاپیلاری, تومورهای خوش‌خیم عروقی
% =========================================================================

% =========================================================================
% CLINICAL EVIDENCE LEDGER & PMIDS (WHO 5th Ed 2024 / Landmark Trials)
% =========================================================================
% 1. WHO Classification of Head and Neck Tumours, 5th Ed (2024): Vascular Anomalies & Tumours.
% 2. Neville BW et al. Oral and Maxillofacial Pathology, 5th Ed (2024): Soft Tissue Neoplasms.
% 3. PMID: 25693013 - Léauté-Labrèze C et al. (N Engl J Med 2015) A randomized, controlled trial of oral propranolol in infantile hemangioma.
% 4. PMID: 30584062 - Krowchuk DP et al. (Pediatrics 2019) Clinical Practice Guideline for the Management of Infantile Hemangiomas.
% 5. PMID: 41446713 - Werpachowski N et al. (J Clin Aesthet Dermatol 2025) Topical Timolol in Infantile Hemangiomas: Systematic Review.
% 6. PMID: 42414159 - da Silva LOP et al. (Oral Surg Oral Med Oral Pathol Oral Radiol 2026) Ultrasonographic Parameters & Sclerotherapy of Oral Vascular Anomalies.
% 7. PMID: 35468033 - De Ravin E et al. (Pediatr Hematol Oncol 2022) International Clinical Practice Guidelines on Hemangioma: Systematic Review.
% =========================================================================

\begin{document}

\begin{center}
{\Large\bfseries فلوچارت تصمیم‌گیری بالینی: همانژیوم کاپیلاری دهان (\lr{Capillary Hemangioma})}\\[2pt]
{\small\itshape تریاژ دیاسکوپی و همودینامیک، ایمونوهیستوشیمی \lr{GLUT-1}، پروتکل بتابلوکر سیستمیک/موضعی، اسکلروتراپی و کنترل خونریزی}
\end{center}
\vspace{-4pt}
\hrule height 0.8pt
\vspace{6pt}

\begin{center}
\begin{latin}
\linespread{0.95}\selectfont
\begin{tikzpicture}[
    every node/.style={font=\sffamily\footnotesize, align=center},
    stepbox/.style={rectangle, draw=black, fill=white, line width=0.7pt, inner sep=4.5pt},
    decision/.style={diamond, aspect=2.0, draw=black, fill=white, line width=0.7pt, inner sep=2.5pt, font=\sffamily\scriptsize\bfseries},
    arrow/.style={-{Stealth[scale=0.88]}, thick, black}
]

% 1. Initial Presentation
\node (n1) [stepbox, line width=1.0pt, text width=12.2cm] at (0, 0) {
    \rl{\textbf{تظاهر بالینی: ندول/ماکول عروقی قرمز تا ارغوانی مخاط دهان (\lr{Oral Vascular Lesion})}}\\[2pt]
    \scriptsize \rl{تست دیاسکوپی مثبت (رنگ‌پریدگی با فشار \lr{Blanching}) $\bullet$ زبان، لب، مخاط باکال یا لثه $\bullet$ رد تپش یا سوفل (\lr{Bruit/Thrill})}
};

% 2. Mandatory Hemodynamic Triage
\node (n2) [stepbox, text width=14.2cm] at (0, -1.85) {
    \colorbox{hlyellow}{\strut\textbf{\rl{مرحله ۱ الزامی: تریاژ همودینامیک و داپلر (رد قطعی مالفورماسیون شریانی-وریدی پرجریان)}}}\\[2pt]
    \scriptsize \rl{سونوگرافی کالر داپلر (\lr{CDUS}) جهت ارزیابی سرعت جریان $\bullet$ رادیوگرافی/\lr{CBCT} جهت رد درگیری داخل استخوانی فکین}
};

% 3. High-Flow / Intrabony Decision
\node (d1) [decision] at (0, -3.85) {\rl{جریان بالا / استخوان؟}\\[-1pt]\scriptsize \lr{High-Flow / Bone?}};
\node (r_exit) [stepbox, draw=darkred, line width=1.0pt, text width=3.8cm] at (5.8, -3.85) {
    \rl{\textbf{مالفورماسیون پرجریان (\lr{AVM}) \emoji{⚠}}}\\[2pt]
    \scriptsize \colorbox{hlred}{\strut\color{darkred}\textbf{\rl{بیوپسی اکیداً ممنوع}}}\\[2pt]
    \scriptsize \rl{ارجاع اورژانسی به آنژیوگرافی}\\
    \rl{و آمبولیزاسیون مداخله‌ای}
};

% 4. Mandatory Biologic Characterization
\node (n3) [stepbox, line width=1.0pt, text width=15.8cm] at (0, -6.15) {
    \colorbox{hlyellow}{\strut\textbf{\rl{ارزیابی بیولوژیک و ایمونوهیستوشیمی طبقه‌بندی \lr{ISSVA / WHO 5th Ed 2024}}}}\\[2pt]
    \scriptsize \rl{تمایز تومور حقیقی از مالفورماسیون عروقی: مارکر \lr{GLUT-1} (در اینفانتایل مثبت، مالفورماسیون منفی) $\bullet$ فاز پرولیفراتیو در برابر اینولوشن}
};

% 5. Clinical Phenotype Decision
\node (d2) [decision] at (0, -8.35) {\rl{فنوتیپ بالینی}\\[-1pt]\scriptsize \lr{Clinical Phenotype?}};

% 6. 4 Main Clinical Presentation Columns
\node (h1) [stepbox, text width=3.6cm] at (-6.15, -10.65) {
    \rl{\textbf{سطحی / نوزادی کم‌خطر}}\\[2pt]
    \scriptsize \rl{ضایعه مخاطی خفیف، کوچکتر از ۲ سانتی‌متر}\\\rl{فاقد اولسر و اختلال بلع یا تنفس}\\\rl{شانس پس‌رفت: }\colorbox{hlgray}{\strut\textbf{\rl{تا ۹۰٪ تا ۹ سالگی}}}
};
\node (h2) [stepbox, text width=3.6cm] at (-2.05, -10.65) {
    \rl{\textbf{لوبولار مخاطی (\lr{LCH/PG})}}\\[2pt]
    \scriptsize \rl{ندول لب یا مخاط باکال ناشی از تروما}\\\rl{خونریزی‌های مکرر با غذا خوردن}\\\rl{رشد سریع اگزوفیتیک در اطفال و جوانان}
};
\node (h3) [stepbox, line width=0.9pt, draw=orange!80!black, text width=3.6cm] at (2.05, -10.65) {
    \rl{\textbf{پرولیفراتیو شدید / اولسره}}\\[2pt]
    \scriptsize \rl{رشد سریع مخاطی، اختلال تغذیه}\\\rl{زخم دردناک یا تهدید راه هوایی فک}\\\rl{اندیکاسیون: }\colorbox{hlred}{\strut\color{darkred}\textbf{\rl{شروع فوری درمان دارویی}}}
};
\node (h4) [stepbox, draw=darkred, line width=1.2pt, text width=3.6cm] at (6.15, -10.65) {
    \rl{\textbf{کونژنتال پایدار (\lr{NICH}) / مالفورماسیون}}\\[2pt]
    \scriptsize \rl{ضایعه عروقی غیرپس‌رونده در بدو تولد}\\\rl{مارکر \lr{GLUT-1 (-)}؛ مقاومت به بتابلوکر}\\\rl{توده عروقی عمقی و آسیمتری بافت نرم}
};

% 7. Primary Interventions
\node (i1) [stepbox, text width=3.6cm] at (-6.15, -13.15) {
    \rl{\textbf{تیمولول موضعی / پایش فعال}}\\[2pt]
    \scriptsize \lr{Timolol 0.5\% GFS} \rl{۲ قطره دوبار در روز}\\\rl{یا پیگیری بدون مداخله تهاجمی}\\\lr{[PMID: 41446713 - 2025]}
};
\node (i2) [stepbox, text width=3.6cm] at (-2.05, -13.15) {
    \rl{\textbf{اکسیزیون جراحی یا لیزر}}\\[2pt]
    \scriptsize \rl{برداشت کامل تا لایه ساب‌موکوزا}\\\rl{هموستاز با لیزر دیود یا کوتر}\\\rl{میزان عود: }\colorbox{hlgray}{\strut\textbf{\rl{کمتر از ۳٪}}}
};
\node (i3) [stepbox, line width=0.9pt, draw=orange!80!black, text width=3.6cm] at (2.05, -13.15) {
    \rl{\textbf{پروپرانولول خوراکی خط اول \emoji{🎯}}}\\[2pt]
    \scriptsize \lr{Propranolol 2--3\,mg/kg/day}\\\rl{منقسم در دو دوز همراه با شیر/تغذیه}\\\lr{[PMID: 25693013 - NEJM RCT]}
};
\node (i4) [stepbox, draw=darkred, line width=1.2pt, text width=3.6cm] at (6.15, -13.15) {
    \rl{\textbf{اسکلروتراپی هدایت‌شده}}\\[2pt]
    \scriptsize \rl{تزریق فوم پولیدوکانول ۱٪ تا ۲٪}\\\rl{یا بلئومایسین تحت گاید سونوگرافی}\\\lr{[PMID: 42414159 - da Silva 2026]}
};

% 8. Recall Schedules
\node (s1) [stepbox, text width=3.6cm] at (-6.15, -15.65) {
    \rl{\textbf{زمان‌بندی پیگیری:}}\\[2pt]
    \colorbox{hlgray}{\strut\textbf{\rl{هر ۶ ماه یکبار}}}\\[2pt]
    \scriptsize \rl{فتوگرافی سریال و پایش پس‌رفت}
};
\node (s2) [stepbox, text width=3.6cm] at (-2.05, -15.65) {
    \rl{\textbf{زمان‌بندی پیگیری:}}\\[2pt]
    \colorbox{hlgray}{\strut\textbf{\rl{ماه ۱ و ۳ پس از عمل}}}\\[2pt]
    \scriptsize \rl{بررسی هیلینگ مخاط و پاتولوژی}
};
\node (s3) [stepbox, line width=0.9pt, draw=orange!80!black, text width=3.6cm] at (2.05, -15.65) {
    \rl{\textbf{زمان‌بندی پیگیری:}}\\[2pt]
    \colorbox{hlred}{\strut\color{darkred}\textbf{\rl{ماهانه در فاز رشد}}}\\[2pt]
    \scriptsize \rl{تنظیم دوز پروپرانولول با وزن کودک}
};
\node (s4) [stepbox, draw=darkred, line width=1.2pt, text width=3.6cm] at (6.15, -15.65) {
    \rl{\textbf{زمان‌بندی پیگیری:}}\\[2pt]
    \colorbox{hlred}{\strut\color{darkred}\textbf{\rl{هر ۳ تا ۶ ماه}}}\\[2pt]
    \scriptsize \rl{سونوگرافی داپلر و پایش ترومبوز}
};

% 9. Trigger Box
\node (trigger) [stepbox, line width=0.9pt, text width=16.4cm, align=right] at (0, -18.35) {
    \centerline{\rl{\textbf{علائم خطر بالینی (\lr{Red Flags}) و اندیکاسیون‌های بستری و مداخله اورژانسی:}}}\par\vspace{3pt}
    \scriptsize \rl{\colorbox{hlyellow}{\strut\textbf{مداخله فوری الزامی است در:}} (۱) استریدور یا انسداد حاد راه هوایی ناشی از گسترش ضایعه به زبان یا حلق $\bullet$ (۲) خونریزی فعال شدید غیرقابل مهار ناشی از ترومای مخاطی $\bullet$ (۳) اولسراسیون وسیع همراه با درد شدید و ناتوانی در تغذیه نوزاد $\bullet$ (۴) عوارض بتابلوکر سیستمیک شامل برادیکاردی (کمتر از ۸۰ ضربان در دقیقه)، هیپوگلیسمی علامت‌دار یا اسپاسم برونش $\bullet$ (۵) ضایعات پلاک‌مانند وسیع صورت با شک به سندرم \lr{PHACE}.}
};

% 10. Refractory Box
\node (refract) [stepbox, text width=16.4cm, align=right] at (0, -20.85) {
    \centerline{\rl{\textbf{پروتکل پله‌ای برای موارد مقاوم به درمان دارویی / ضایعات عروقی پیچیده:}}}\par\vspace{3pt}
    \scriptsize \rl{\textbf{سطح ۳ (مقاوم به پروپرانولول):} پردنیزولون خوراکی (\lr{2--3\,mg/kg/day}) یا اضافه کردن لیزر رنگی پالسی (\lr{PDL 595\,nm}) جهت ترمیم اولسر و محو عروق سطحی.}\par\vspace{2.5pt}
    \scriptsize \rl{\textbf{سطح ۴ (مهارکننده تارگت بیولوژیک \lr{mTOR}):} سیرولیموس خوراکی (\lr{Sirolimus 0.8\,mg/m$^2$ q12h} با هدف سطح خونی \lr{5--10\,ng/mL}) در ضایعات مالفورماسیون عروقی و لنفاتیک مقاوم.}\par\vspace{2.5pt}
    \scriptsize \rl{\textbf{سطح ۵ (سالواژ مداخله‌ای و جراحی):} آمبولیزاسیون ترانس‌آرتریال سوپرسلکتیو به همراه اکسیزیون جراحی کامل تحت بیهوشی با کوتر بای‌پولار یا تیغ اولتراسونیک هارمونیک.}
};

% Connections
\draw [arrow] (n1) -- (n2);
\draw [arrow] (n2) -- (d1);
\draw [arrow] (d1.east) -- node[above, font=\scriptsize\bfseries] {\rl{بله (خطرناک)}} (r_exit.west);
\draw [arrow] (d1.south) -- node[right=4pt, font=\scriptsize\bfseries] {\rl{خیر (کم‌جریان/مخاطی)}} (n3.north);
\draw [arrow] (n3.south) -- (d2.north);

% Bus-Bar from Decision 2 to 4 Histology columns
\draw [thick, black] (d2.south) -- (0, -9.45);
\draw [thick, black] (-6.15, -9.45) -- (6.15, -9.45);
\draw [arrow] (-6.15, -9.45) -- (h1.north);
\draw [arrow] (-2.05, -9.45) -- (h2.north);
\draw [arrow] (2.05, -9.45) -- (h3.north);
\draw [arrow] (6.15, -9.45) -- (h4.north);

% Column Downward Arrows
\draw [arrow] (h1.south) -- (i1.north);
\draw [arrow] (h2.south) -- (i2.north);
\draw [arrow] (h3.south) -- (i3.north);
\draw [arrow] (h4.south) -- (i4.north);

\draw [arrow] (i1.south) -- (s1.north);
\draw [arrow] (i2.south) -- (s2.north);
\draw [arrow] (i3.south) -- (s3.north);
\draw [arrow] (i4.south) -- (s4.north);

% Collector Bus into Trigger Box
\draw [thick, black] (s1.south) -- (-6.15, -16.95);
\draw [thick, black] (s2.south) -- (-2.05, -16.95);
\draw [thick, black] (s3.south) -- (2.05, -16.95);
\draw [thick, black] (s4.south) -- (6.15, -16.95);
\draw [thick, black] (-6.15, -16.95) -- (6.15, -16.95);
\draw [arrow] (0, -16.95) -- (trigger.north);

\draw [arrow] (trigger.south) -- (refract.north);

\end{tikzpicture}
\end{latin}
\end{center}

\end{document}
